\documentclass{article}

\usepackage[T1]{fontenc}

\usepackage{booktabs,graphicx,siunitx}

\begin{document}

\section{Chi-square independence}

\begin{tabular}{llll}
\toprule
Row \textbackslash{} Column & No & Yes & Total \\
\midrule
Blended & 30 [29.00] (37.5\% / 34.5\%) r=0.28 & 50 [51.00] (62.5\% / 32.7\%) r=$-$0.28 & 80 \\
Flipped & 17 [29.00] (21.3\% / 19.5\%) r=$-$3.42 & 63 [51.00] (78.8\% / 41.2\%) r=3.42 & 80 \\
Lecture & 40 [29.00] (50.0\% / 46.0\%) r=3.13 & 40 [51.00] (50.0\% / 26.1\%) r=$-$3.13 & 80 \\
Total & 87 & 153 & 240 \\
\bottomrule
\end{tabular}

\begin{tabular}{llll}
\toprule
$\chi$\textsuperscript{2} & df & p & Cram\'{e}r's V [95\% CI] \\
\midrule
14.39 & 2 & $<$.001 & 0.24 [0.13, 1.00] \\
\bottomrule
\end{tabular}

\includegraphics{figures/01_chi-square-independence/figure_bar.png}

Tests whether two categorical variables are associated. A p below alpha means they are related; Cram\'{e}r's V gives the strength of association (0--1); each cell's standardized residual (r) beyond about $\pm$1.96 flags a category combination that significantly drives the association. Valid only when expected counts are mostly $\ge$ 5 (the app flags this and suggests Fisher's exact) and each case appears once (raw counts, not percentages). Your significance threshold ($\alpha$) is 0.05.

A chi-square test of independence gave $\chi$\textsuperscript{2}(2, N=240)=14.39, p $<$ .001, V=.24 [.13, 1.00].

Statistical basis: Chi-square test of independence (Pearson, K. (1900). "On the criterion that a given system of deviations from the probable in the case of a correlated system of variables is such that it can be reasonably supposed to have arisen from random sampling." Philosophical Magazine, 50(302), 157-175.); Cram\'{e}r's V effect size (w-family; Cohen, 1988) (Cohen, J. (1988). Statistical Power Analysis for the Behavioral Sciences (2nd ed.). Routledge.).

\end{document}